\documentclass[11pt, a4paper]{report}

% --- Encodage et Langue ---
\usepackage[utf8]{inputenc}
\usepackage[T1]{fontenc}
\usepackage[french]{babel}
\usepackage{hyperref}
\usepackage{markdown}

% --- Mise en page ---
\usepackage[margin=2.5cm]{geometry} % Marges équilibrées
\usepackage{titlesec}
% --- Mathématiques ---
\usepackage{amsmath, amssymb, amsthm}

\usepackage{tabularx}
\usepackage{booktabs}
\usepackage[table]{xcolor}
\usepackage{float}
\definecolor{headerblue}{RGB}{31, 73, 125}

\usepackage{array,multirow,makecell}
\setcellgapes{1pt}
\makegapedcells
\newcolumntype{R}[1]{>{\raggedleft\arraybackslash }b{#1}}
\newcolumntype{L}[1]{>{\raggedright\arraybackslash }b{#1}}
\newcolumntype{C}[1]{>{\centering\arraybackslash }b{#1}}



\title{Compte rendu hebdomadaire \\ Master MIX Première Année}
\author{Matthieu Durand}
\date{}

\begin{document}
\maketitle
\tableofcontents
\newpage

\section*{Tuteurs}
\begin{itemize}
    \item M. Essoham Ali, Maître de Conférence à l'Institut de Mathématiques Appliquées d'Angers
    \item M. Christophe Saint-Jean, Professeur à La Rochelle Université
\end{itemize}

%%%%%%%%%%%%%%%%%%%%%%%%
\section{Information}
Ce document fait office de rapport hebdomadaire, permettant de suivre l'avancement du stage 
et des différentes étapes. 
Dans chaque section, nous décrivons les éléments de la semaine ainsi que les démarche 
suivies. 
%%%%%%%%%%%%%%%%%%%%%%%%

\section*{Introduction}
Ce stage s'est déroulé à l'Institut de Mathématiques Appliquées d'Angers, 
encadré par M. Essoham Ali. L'objectif était de mettre en application des modèles 
étudiés durant mon année de Master afin de prédire l'issue de matchs de football. 
Plusieurs étapes ont été nécessaires : la récupération des données, le nettoyage 
de ces données, et la création de descripteurs fiables évitant le 
\textit{data-leakage}\footnote{wikipédia : In statistics and machine 
learning, leakage (also known as data leakage or target leakage) refers to the use of 
information during model training that would not be available at prediction time.}.



%%%%%%%%%%%%%%%%%%%%%%%%%%%%%%%%%%%
\chapter{Semaine 1}
\section{État de l'art et collecte de données}

% Dans un premier temps, mon travail a consisté à réaliser une revue de l'état de l'art 
% concernant les modèles de prédiction des matchs de football. Pour ce faire, je me suis 
% principalement concentré sur trois publications de référence :
% \begin{itemize}
% \item L'article de \textbf{Dixon et Coles (1997)} \cite{dixon} aborde l'approche 
% statistique du problème via l'implémentation de modèles basés sur la loi de Poisson.
% \item L'article d'\textbf{Egidi et Torelli (2021)} \cite{egidi} poursuit cette analyse en différenciant deux types de modélisations : celles basées sur l'issue du match (\textit{result-based}) et celles basées sur le score exact (\textit{goal-based}).
% \item Enfin, l'article de \textbf{Carloni et al. (2022)} \cite{carloni} explore 
% l'implémentation d'algorithmes de Machine Learning (ML), tels que les réseaux de neurones, pour traiter ce problème.
% \end{itemize}  

% Par la suite, j'ai procédé à la collecte des données relatives aux cinq championnats 
% majeurs européens sur les dix dernières saisons via le site \textit{football-data.co.uk}.
\section{Revue de la littérature}
La phase initiale de ce stage a été consacrée à une étude approfondie de l'état 
de l'art concernant la modélisation prédictive dans le football. Cette revue de 
littérature s'est articulée autour de trois travaux fondateurs qui illustrent 
l'évolution des méthodologies dans ce domaine :

\begin{itemize}
    \item \textbf{Approche statistique classique :} L'étude de 
    \textbf{Dixon et Coles (1997)} \cite{dixon} constitue la référence pour 
    l'application de la loi de Poisson. Les auteurs y introduisent des paramètres 
    d'attaque et de défense spécifiques à chaque équipe pour estimer la probabilité 
    des scores.
    \item \textbf{Comparaison des paradigmes :} \textbf{Egidi et Torelli (2021)} 
    \cite{egidi} enrichissent cette analyse en comparant deux approches distinctes :
     la modélisation directe de l'issue du match (\textit{result-based}) et celle 
     fondée sur le score exact (\textit{goal-based}).
    \item \textbf{Apprentissage automatique :} Enfin, \textbf{Carloni et al. (2021)}
     \cite{carloni} explorent le potentiel des 
    algorithmes d'apprentissage supervisé, notamment les réseaux de neurones, 
    pour capturer des relations non-linéaires complexes que les modèles statistiques 
    traditionnels peinent parfois à isoler.
\end{itemize}

\section{Constitution de la base de données}
Une fois le cadre théorique établi, nous avons procédé à la constitution d'un dataset 
historique. Les données ont été extraites du portail \textit{football-data.co.uk}, 
couvrant les dix dernières saisons des cinq grands championnats européens (Premier 
League, LaLiga, Serie A, Bundesliga et Ligue 1). Cette profondeur historique est 
essentielle pour garantir la robustesse des modèles d'apprentissage que nous 
souhaitons implémenter.

%%%%%%%%%%%%%%%%%%%%%%%%%%%%%%%%%%%
\section{Semaine 2 : Nettoyage et calcul de nouveaux déscripteurs}

Durant cette semaine, l'objectif était de débuter le traitement informatique des
données. Pour cela, j'ai utilisé le langage de programmation Python et
majoritairement la bibliothèque Pandas. Après avoir collecté les données sur
10 saisons et 5 championnats, la première étape a été de nettoyer les données
manquantes et de sélectionner les colonnes pertinentes du dataset.
J'ai notamment supprimer toutes les données relatives aux cotes des paris. En effet, 
ces colonnes sont supprimées pour éviter le Data-Leakage. 
Elles reflètent des probabilités de sortie qui fausseraient l'apprentissage réel du modèle.

Par la suite, j'ai construit de nouvelles caractéristiques à partir des données

brutes. En premier lieu, j'ai calculé le nombre de \textbf{jours de repos} séparant
le match à la date $t$ du match précédent à la date $t-1$. Ce descripteur peut
s'avérer peu pertinent pour les équipes ayant été reléguées puis promues, car
l'interruption entre deux saisons dans des championnats différents introduit un
biais dans le calcul.

J'ai également mis en place des indicateurs de \textbf{forme récente}, en calculant
pour chaque équipe la moyenne de buts marqués sur les 5 derniers matchs. De la même
manière, un indicateur de \textbf{force offensive} a été construit en moyennant le
nombre de tirs tentés sur les 5 derniers matchs, et un indicateur de
\textbf{force défensive} en moyennant le nombre de tirs subis sur la même fenêtre
glissante. Ces trois indicateurs ont été calculés aussi bien pour les statistiques
de fin de match (\textit{Full Time}) que pour les statistiques de mi-temps
(\textit{Half Time}).

Pour le \textbf{score Elo}\footnote{Implémentation inspirée de
\url{https://www.eloratings.net/about}}, chaque équipe part d'un score initial de
$1500$. Ce score est mis à jour après chaque match en fonction de l'issue du match
et de l'écart de buts, via un facteur multiplicatif $G$ qui amplifie la mise à jour
pour les victoires avec un grand écart. Deux scores Elo indépendants ont été
calculés : l'un basé sur les résultats finaux (\textit{FT\_Elo}), l'autre basé sur
les résultats à la mi-temps (\textit{HT\_Elo}), afin de capturer la capacité d'une
équipe à bien débuter les matchs.

Enfin, en s'inspirant de la démarche de Dixon \& Coles --- dont le modèle repose
sur le rapport $\alpha_i / \alpha_j$ entre les forces d'attaque --- et d'Egidi \&
Torelli --- qui construisent explicitement la différence de classement $w_n =
\text{rank}_H - \text{rank}_A$ comme prédicteur --- nous avons ajouté pour chaque
paire de descripteurs une colonne \textbf{\_diff} (différence $H - A$) et une
colonne \textbf{\_ratio} (rapport $H / A$). Ces nouvelles variables permettent au
modèle de disposer directement d'une mesure de l'écart de niveau entre les deux
équipes, sans avoir à l'inférer implicitement à partir des valeurs brutes.


\subsection*{Descriptif des features construites}

L'ensemble des caractéristiques construites est organisé en six groupes fonctionnels.

\begin{itemize}

    \item \textbf{Forme récente} --- \texttt{FT\_Hforme}, \texttt{FT\_Aforme},
    \texttt{HT\_Hforme}, \texttt{HT\_Aforme} : moyenne des buts marqués sur les
    5 derniers matchs, respectivement pour l'équipe à domicile (H) et l'équipe
    visiteuse (A). Cet indicateur est calculé aussi bien sur les statistiques de
    fin de match (\textit{Full Time}) qu'à la mi-temps (\textit{Half Time}).

    \item \textbf{Force offensive} --- \texttt{FT\_Hatt}, \texttt{FT\_Aatt} :
    moyenne du nombre de tirs tentés sur les 5 derniers matchs. Proxy de la
    pression offensive exercée par l'équipe.

    \item \textbf{Force défensive} --- \texttt{FT\_Hdef}, \texttt{FT\_Adef},
    \texttt{HT\_Hdef}, \texttt{HT\_Adef} : moyenne du nombre de tirs subis sur
    les 5 derniers matchs (\textit{Full Time}) et des buts concédés à la
    mi-temps (\textit{Half Time}). Plus cette valeur est faible, plus la défense
    de l'équipe est solide.

    \item \textbf{Précision} --- \texttt{FT\_Hprecision}, \texttt{FT\_Aprecision} :
    rapport entre le nombre moyen de tirs cadrés et le nombre moyen de tirs tentés
    sur les 5 derniers matchs. Mesure la qualité des occasions créées.
    La variante pondérée \texttt{FT\_Hprec\_weight}, \texttt{FT\_Aprec\_weight}
    donne un poids supplémentaire (facteur $w_g = 1{,}5$) aux tirs convertis en
    buts, afin de valoriser l'efficacité offensive.

    \item \textbf{Repos} --- \texttt{HRepos}, \texttt{ARepos} : nombre de jours
    écoulés depuis le dernier match joué par l'équipe (domicile ou extérieur),
    plafonné à 25 jours pour limiter les biais liés aux promotions et relégations.

    \item \textbf{Score Elo} --- \texttt{FT\_Elo\_H}, \texttt{FT\_Elo\_A},
    \texttt{HT\_Elo\_H}, \texttt{HT\_Elo\_A} : score Elo de chaque équipe avant
    le match, calculé itérativement dans l'ordre chronologique avec un score
    initial de $1500$. La mise à jour tient compte de l'issue et de l'écart de
    buts via le facteur $G$. Un score Elo indépendant est calculé sur les
    résultats à la mi-temps afin de capturer la tendance d'une équipe à bien ou
    mal démarrer les matchs.

    \item \textbf{Fautes et cartons} --- \texttt{FT\_HavgF}, \texttt{FT\_AavgF},
    \texttt{FT\_HavgY}, \texttt{FT\_AavgY}, \texttt{FT\_HavgR}, \texttt{FT\_AavgR} :
    moyennes respectivement du nombre de fautes commises, de cartons jaunes et de
    cartons rouges reçus sur les 5 derniers matchs. Ces indicateurs renseignent sur
    le style de jeu de l'équipe et sa discipline, les cartons rouges étant
    potentiellement corrélés aux matchs joués en infériorité numérique.

\end{itemize}

Pour chacun des groupes ci-dessus, trois colonnes dérivées sont systématiquement
ajoutées afin de faciliter la tâche des modèles de classification :

\begin{itemize}
    \item \textbf{\_diff} : différence $H - A$ entre la valeur de l'équipe à
    domicile et celle de l'équipe visiteuse. Par exemple,
    \texttt{FT\_forme\_diff} $=$ \texttt{FT\_Hforme} $-$ \texttt{FT\_Aforme}.
    Cette approche s'inspire directement d'Egidi \& Torelli, qui construisent
    $w_n = \text{rank}_H - \text{rank}_A$ comme prédicteur explicite de l'écart
    de niveau.

    \item \textbf{\_ratio} : rapport $H / A$ entre les deux valeurs. Cette
    formulation est cohérente avec la structure multiplicative du modèle de
    Dixon \& Coles, dans lequel la probabilité d'issue dépend du rapport
    $\alpha_i / \alpha_j$ entre les forces d'attaque des deux équipes.

    \item \textbf{\_dif} (pour l'Elo uniquement) : déjà inclus sous la forme
    \texttt{FT\_Elo\_dif} $=$ \texttt{FT\_Elo\_H} $-$ \texttt{FT\_Elo\_A} et
    \texttt{HT\_Elo\_dif}.
\end{itemize}

Au total, \textbf{53 features} sont ainsi construites et sauvegardées dans un fichier csv
, auxquelles s'ajoutent les trois colonnes cibles
\texttt{Hvs}, \texttt{Avs} et \texttt{Dvs} pour les analyses
\textit{One-vs-All}.

Trois colonnes cibles ont également été ajoutées --- \texttt{Hvs}, \texttt{Avs} et
\texttt{Dvs} --- afin de permettre une analyse selon le paradigme
\textit{One-vs-All} couramment utilisé en apprentissage automatique pour les
problèmes de classification multi-classes.

%%%%%%%%%%%%%%%%%%%%%%%%%%%%%%%%%%%
\chapter{Semaine 3 }

L'objectif de cette est de poursuivre l'analyse et le traitement de la 
base de données. 
Par la suite nous allons débuter l'Implémentation de modèle de classification 
afin de débuter les prédictions des modèles.


\section{Séléction des features}

Un fois les déscripteurs calculés, nous devons choisir les quels conserver pour obtenir 
les meilleurs résultats possible. 
Dans cet objectif nous avons explorer la méthode RFE(recursive feature elimination). 
Étant donné un estimateur externe qui pondère les caractéristiques, l'objectif 
de l'élimination récursive de caractéristiques (RFE) est de sélectionner les 
caractéristiques en considérant successivement des ensembles de plus en plus petits.
La variant que nous avons utilisé est RFECV, CV signifie simplement validation 
croisée (Cross validation). Le nombre de caractéristiques sélectionnées est ajusté 
automatiquement en adaptant un sélecteur RFE aux différentes plis de validation 
croisée. Dans notre cas il est important de ne pas oublier que les données 
possèdent un ordre chronologique, ainsi nous devons utiliser un \textit{train/test split} 
adéquat, c'est pour cela que nous avons utilisé \textit{TimeSeriesSplit}.  

%%%%%%%%%%%%%%%%%%%%%%%%%%%%%%%%%%%
% \input{s4/s4}

%%%%%%%%%%%%%%%%%%%%%%%%%%%%%%%%%%%
% \input{s5/s5}

%%%%%%%%%%%%%%%%%%%%%%%%%%%%%%%%%%%
% \input{s6/s6}

%%%%%%%%%%%%%%%%%%%%%%%%%%%%%%%%%%%
% \input{s7/s7}

%%%%%%%%%%%%%%%%%%%%%%%%%%%%%%%%%%%
% \input{s8/s8}

%%%%%%%%%%%%%%%%%%%%%%%%%%%%%%%%%%%
%%%%%%%%%%%%%%%%%%%%%%%%%%%%%%%%%%%
\chapter{Modèles}
\section{Listes des modèles à tester}
mèf dans notre travail, nous cherchons à prédire l'issue non 
pas le score finale. Cette liste est amenée à évoluer. 
De même pour pour les métriques.
\begin{itemize}
    \item Regression logistique
    \item Random Forest
    \item KNN
    \item SVM
    \item XGBoosting
    \item CatBoosting
    \item RNN (LSTM)
    % \item 
\end{itemize}
\section{Listes des métriques d'évaluation}


\begin{itemize}
    \item ROI
    \item f1-score
    \item log-loss
    \item brier-score
    \item roc auc 
    \item calibration curve
\end{itemize}
%%%%%%%%%%%%%%%%%%%%%%%%%%%%%%%%%%%
\newpage
\chapter*{Annexe}

\section{Documentation des Données Football}

Ce document détaille la structure des données utilisées pour l'analyse et la 
prédiction de matchs de football. Les colonnes sont regroupées par catégories 
thématiques.

\subsection{Informations Générales \& Identification}

Données de base permettant d'identifier le contexte du match.

\begin{table}[h]
\centering
\begin{tabular}{|l|l|}
\hline
\textbf{Colonne} & \textbf{Description} \\ \hline
\texttt{Div} & Division (championnat) \\ \hline
\texttt{Season} & Saison \\ \hline
\texttt{League} & Nom de la ligue (ex : Premier League) \\ \hline
\texttt{Date} & Date du match (aaaa/mm/jj) \\ \hline
\texttt{HomeTeam} & Équipe évoluant à domicile \\ \hline
\texttt{AwayTeam} & Équipe évoluant à l'extérieur \\ \hline
\end{tabular}
\end{table}

\subsection{Résultats \& Statistiques de Match (Brut)}

Données constatées à la fin de la rencontre.

\subsubsection{Fin de match (Full Time)}
\begin{itemize}
    \item \textbf{FTHG / FTAG} : Buts marqués par l'équipe à domicile / extérieur.
    \item \textbf{FTR} : Résultat final (\textbf{1} : Domicile, \textbf{2} : 
    Extérieur, \textbf{0} : Nul).
\end{itemize}

\subsubsection{Mi-temps (Half Time)}
\begin{itemize}
    \item \textbf{HTHG / HTAG} : Buts à la mi-temps (Dom/Ext).
    \item \textbf{HTR} : Résultat à la mi-temps 
    (\textbf{1} : Domicile, \textbf{2} : Extérieur, \textbf{0} : Nul).
\end{itemize}

\subsubsection{Statistiques de Jeu}
\begin{itemize}
    \item \textbf{HS / AS} : Tirs (Dom/Ext).
    \item \textbf{HST / AST} : Tirs cadrés (Dom/Ext).
    \item \textbf{HC / AC} : Corners (Dom/Ext).
    \item \textbf{HF / AF} : Fautes commises (Dom/Ext).
    \item \textbf{HY / AY} : Cartons jaunes (Dom/Ext).
    \item \textbf{HR / AR} : Cartons rouges (Dom/Ext).
\end{itemize}

\subsection{Indicateurs de Performance (Calculés)}

Moyennes mobiles et indicateurs d'efficacité basés sur les matchs précédents 
($XX \in \{FT, HT\}$).

\subsubsection{Domicile (Home) vs Extérieur (Away)}

\begin{table}[h]
\centering
\resizebox{\textwidth}{!}{
\begin{tabular}{|l|l|l|l|}
\hline
\textbf{Catégorie} & \textbf{Domicile} & \textbf{Extérieur} & \textbf{Description} \\ \hline
\textbf{Forme (FT/HT)} & \texttt{XX\_Hforme} & \texttt{XX\_Aforme} & Moyenne des 
buts marqués (5 derniers matchs) \\ \hline
\textbf{Attaque (FT)} & \texttt{FT\_Hatt} & \texttt{FT\_Aatt} & Moyenne des tirs 
effectués \\ \hline
\textbf{Défense (FT/HT)} & \texttt{XX\_Hdef} & \texttt{XX\_Adef} & Moyenne des 
tirs concédés \\ \hline
\textbf{Elo Score (FT/HT)} & \texttt{XX\_Elo\_H} & \texttt{XX\_Elo\_A} & Niveau 
de puissance avant match \\ \hline
\textbf{Précision} & \texttt{FT\_Hprecision} & \texttt{FT\_Aprecision} & Ratio 
Tirs cadrés / Tirs totaux \\ \hline
\textbf{Efficacité} & \texttt{FT\_Hprec\_weight} & \texttt{FT\_Aprec\_weight} & 
Précision pondérée (poids buts = 1.5) \\ \hline
\end{tabular}
}
\end{table}

\subsubsection{Discipline \& Repos}

\begin{table}[h]
\centering
\begin{tabular}{|l|l|l|l|}
\hline
\textbf{Catégorie} & \textbf{Dom (H)} & \textbf{Ext (A)} & \textbf{Description} \\ \hline
Fautes & \texttt{FT\_HavgF} & \texttt{FT\_AavgF} & Moyenne fautes (5 matchs) \\ \hline
Carton J. & \texttt{FT\_HavgY} & \texttt{FT\_AavgY} & Moyenne jaunes (5 matchs) \\ \hline
Carton R. & \texttt{FT\_HavgR} & \texttt{FT\_AavgR} & Moyenne rouges (5 matchs) \\ \hline
Repos & \texttt{HRepos} & \texttt{ARepos} & Jours depuis dernier match (max 25) \\ \hline
\end{tabular}
\end{table}

\subsection{Variables de Comparaison (Diff \& Ratios)}

Indicateurs pour mettre en évidence l'écart entre les adversaires 
($XX \in \{FT, HT\}$).

\begin{itemize}
    \item \textbf{Formule Différence :} $stat_{home} - stat_{away}$
    \item \textbf{Formule Ratio :} $\frac{stat_{home}}{stat_{away} + 10^{-6}}$
\end{itemize}

\begin{table}[h]
\centering
\begin{tabular}{|l|l|l|}
\hline
\textbf{Catégorie} & \textbf{Différence} & \textbf{Ratio} \\ \hline
Forme & \texttt{XX\_forme\_diff} & \texttt{XX\_forme\_ratio} \\ \hline
Attaque & \texttt{XX\_att\_diff} & \texttt{XX\_att\_ratio} \\ \hline
Défense & \texttt{XX\_def\_diff} & \texttt{XX\_def\_ratio} \\ \hline
Elo & \texttt{XX\_Elo\_diff} & \texttt{XX\_Elo\_ratio} \\ \hline
\end{tabular}
\end{table}

\subsection{Variables Cibles (Labels)}

Colonnes utilisées pour l'entraînement (Machine Learning) : \texttt{Hvs} 
(Victoire D.), \texttt{Avs} (Victoire E.), \texttt{Dvs} (Nul).

\subsection{Cotes des Bookmakers (Supprimées)}

\begin{quote}
    \textbf{Note :} Ces colonnes sont supprimées pour éviter le \textbf{Data Leakage}. 
    Elles reflètent des probabilités de sortie qui fausseraient l'apprentissage réel 
    du modèle.
\end{quote}

%%%%%%%%%%%%%%%%%%%%%%%%%%%%%%%%%%%
\newpage
\nocite{*}
\bibliographystyle{plain}
\bibliography{references}

\end{document}